\newpage
\section{Introduction} 
Empirical research at the intersection of the sociology of culture and stratification has been primarily concerned  with exploring inequalities linked to social position in \emph{access} to and \emph{engagement} with symbolic goods  \citep{katzgerro04}.  The primary question in this burgeoning literature is whether taste and culture consumption patterns---as constitutive of the lifestyles that mark Weberian status groups---continue to be associated with social position \citep{tomlinson03}.  This organizing concern can be appreciated whether the construct of ``social position'' is conceptualized in terms of educational attainment \citep{ warde_etal08}, membership in broadly defined occupational groups \citep{tomlinson03}, aggregate measures of social class \citep{katzgerro06}, or constructed scalar indices of social status based on patterns of interaction (e.g. inter-marriage) across occupational groups \citep{chan_goldthorpe07a}. 

The main conclusion of this still vibrant line of research is that lifestyle practices continue to be strongly shaped---albeit in increasingly complex ways---by social position however it is defined \citep{bennett_etal08}.  One basic finding is now well-established:  members of the educated, salaried-professional elite are more likely to actively appropriate a wider variety of aesthetic offerings, within and across industry-defined genre categories and institutional contexts of dissemination and reception \citep{peterson05, warde08, lizardo08, prieur_savage11}. In essence, what it means to be high-status in contemporary societies is to be an equal opportunity connoisseur of diverse cultural offerings and aesthetic experiences across domains (e.g. literature versus film, and within domains across substyles), forms of delivery (e.g. non-profit versus commercial systems) and external legitimation (e.g. ``prescribed'' or ``authorized'' culture versus ``pop culture'', ``folk culture'' or ``kitsch'') \citep{peterson_kern96, bennett_etal08, prieur_savage11}.
	
This set of findings run counter to the predictions of ``postmodernization'' and class ``destructuration'' theorists who foresaw an inexorable decline in (and ultimate elimination of) the link between social position and lifestyle in market-dominated ``consumer societies''\citep[e.g.][]{pakulski_waters96}. In direct contrast to this claim, culture and stratification scholars can confidently assert that generalized markers of social status and education (rather than social class more narrowly defined) provide the best lens to understand cultural stratification patterns in the current context, confirming the basic findings of an earlier generation of research \citep{dimaggio_useem78a}. 

\subsection{What is good taste good for?}
The persistent focus of the cultural stratification literature on the question of \emph{whether} lifestyle practices continue to be unevenly segmented according to class and status markers has resulted in a relative lack of attention to an equally important question:  what is the \emph{role} (or ``function'') of aesthetic consumption in the everyday lives of members of different social strata? More precisely, \emph{given} the highly educated and those who belong to high-prestige occupational groups are more likely to engage with and be more generally knowledgeable about popular culture and the arts, how is this knowledge implicated in everyday behavior?  In essence, now that we know that eclecticism is the form that ``good'' taste takes in contemporary societies, it is time to ask---with \citep{erickson91}---``what is [this kind of] good taste good for''?

This gap in our knowledge is surprising in light of the fact that the status of popular culture and the arts as a form of \emph{class culture} was, and continues to be, central to the theoretical tradition from which most analysts of the culture-stratification build on in the first place \citep{collins75, dimaggio87, douglas_isherwood96}.  In this respect, Gary Alan Fine's basic complaint about research of popular culture in the early 1970s---that this research generally ignored the \emph{behavioral implications} of everyday engagement with culture and the arts especially as it pertains to face-to-face interaction \citep[453]{fine77}---can be repeated vis a vis research on cultural stratification today.  

One of the key analytic lessons of some of the foundational writings in the culture and stratification literature is that general familiarity with a given class culture is useful symbolic resource when it can be used as a signal providing an entry-point into delimited arenas of association \citep{mohr_dimaggio95}.  This is the reason why, in their influential paper, \citet{lamont_lareau88} linked the notion of cultural capital to an explicit conceptual framework that highlighted the active deployment of a given repertoire of cultural tastes and lifestyle practices for the purposes of drawing symbolic boundaries.  \citet{dimaggio87} was also adamant---building on the early work of Collins \citeyearpar{collins75}---in proposing that the major differences across status groups in the role that engagement with symbolic goods play in everyday social settings is connected to the \emph{routine exploitation of cultural knowledge in social interaction within everyday settings} \citep{erickson96, lizardo06a}. From this perspective, one substantively important way that the set of cultural aptitudes and interests characteristic of upper-middle class lifestyles become actualized as cultural \emph{capital}---that is ``convertible'' into other resources \citep{bourdieu86}---is through the linkage of cultural engagement to everyday processes of informal socialization and thus status group reproduction in the Weberian sense \citep{dimaggio_mohr85}.

